%! TEX program = xelatex
% 2026 MathorCup大数据竞赛初赛模板；其他届次先核官方规范。
% 摘要第一页、目录第二页，正文从1编号；正文≤30页，附录不限。
% 本模板是骨架，文字与人工审核须真实补齐，不能作为已完成论文提交。
\documentclass[UTF8,a4paper,12pt]{ctexart}
\usepackage[left=2.5cm,right=2.5cm,top=2.5cm,bottom=2.5cm]{geometry}
\usepackage{amsmath,amssymb,booktabs,longtable,graphicx,fvextra,xurl}
\usepackage[hidelinks]{hyperref}
\setlength{\headheight}{0pt}
\pagestyle{empty}
\title{论文题目（按本题实际方法与对象填写）}
\author{}
\date{}
\begin{document}
\maketitle
\thispagestyle{empty}
\begin{abstract}
摘要依据实际完成的方法、计算结果和检验填写；关键数值追溯到冻结的结果键。
\par\noindent\textbf{关键词：}按实际对象和方法填写
\end{abstract}
\clearpage
\tableofcontents
\clearpage
\pagestyle{plain}
\pagenumbering{arabic}
\section{问题重述}
逐问列任务、输入、输出、约束、单位、交付物与来源页码。
\section{问题分析}
说明数据可用性、跨问接口与实际可识别范围。
\section{模型假设与符号说明}
区分题给事实、待定解释和团队设计假设。
\section{数据来源与预处理}
记录血缘、时空支持、有效值、样本划分及泄漏检查。
\section{模型建立与求解}
按实际问题数展开；写清数学表达、求解步骤、实际答案和检验。
\section{结果验证与独立复算}
分开同口径数值误差、不同定义的估计量、未覆盖项及实际分歧。
\section{模型评价与结论}
将局部限定同步到摘要、结论和图注。
\begin{thebibliography}{99}
% 仅列实际查阅采用的来源和使用的AI工具，正文对应位置引用。
\end{thebibliography}
\section*{人工智能工具使用声明}
% 本届要求位于参考文献后；正文对应位置标注，程序注释记录辅助编码。
按真实使用范围填写。队员独立核心工作、采纳、人工核验和签署状态不得由代理代填。
\appendix
\section{支撑材料与软件环境}
% support_pack完成后实际引入清单：
% \input{../findings/support_manifest.tex}
列出实际软件版本、运行命令与复现说明。
\section{全部自编源程序}
% 逐个引用真实完整源码，不能仅写文件名或放算法节选：
% \VerbatimInput[breaklines=true,fontsize=\small]{../scripts/solve.py}
\end{document}
